%!TEX root = ../thesis.tex

% A German abstract is mandatory in addition to the English one for a thesis
% written in English. thesis.tex inputs this file inside
% \begin{otherlanguage}{ngerman}, so write plain German here.

\addchap*{Kurzfassung}
\markboth{Kurzfassung}{}

% siunitx has no babel integration; the comma is local to the otherlanguage
% group thesis.tex inputs this file in.
\sisetup{output-decimal-marker={,}}

Coverage-gesteuertes Greybox-Fuzzing verdankt seinen Erfolg beim Finden von Fehlern dem
Durchsatz: Es führt Millionen von Eingaben aus und behält jene, die neuen Code erreichen, ohne je
ein explizites Modell davon zu bilden, was eine gute Eingabe ausmacht.
Wie das Ausführungsbudget verteilt wird, entscheiden handgetunte Heuristiken, und Zielprogramme,
deren Verhalten nicht von einer einzelnen Eingabe, sondern von einer Folge vorheriger Eingaben
abhängt, passen nicht in das zustandslose Modell, das dabei angenommen wird.

Wir modellieren emulatorbasiertes Greybox-Fuzzing als Markov-Entscheidungsprozess und steuern es
mit EfficientZero, einer stichprobeneffizienten MuZero-Variante, die mittels
Monte-Carlo-Baumsuche in einem gelernten latenten Modell ihrer Umgebung plant.
Die unter Emulation erfasste Basisblock-Coverage dient zugleich als Beobachtung und als Belohnung.
Wir formulieren zwei Aktionsräume: Aktionen auf Eingabeebene, bei denen eine Aktion einer Eingabe
entspricht, die das Programm konsumiert, und Aktionen auf Korpusebene, bei denen eine Aktion eine
Position in einem Seed und den dort geschriebenen Bytewert auswählt.
Beide sind in einer wiederverwendbaren Rust-Implementierung von EfficientZero auf Basis des
Deep-Learning-Frameworks Burn und im Fuzzer \emph{mugiwara} umgesetzt, der seine Zielprogramme
unter dem Emulations-Framework MOGI ausführt.

Wir evaluieren sechs Zielprogramme.
Der Agent löst alle drei Ziele mit Aktionen auf Eingabeebene: Er rekonstruiert einen 16~Byte langen
Passcode Position für Position und aktiviert eine OPC~UA-Sitzung gegen open62541.
Bei den Zielen mit Aktionen auf Korpusebene übertrifft er bei identischem Ausführungsbudget eine
ungerichtete Kontrolle und erreicht auf cJSON im Median \num{645} abgedeckte Blöcke gegenüber
\num{542}, im besten von vierzehn Ablationsläufen \num{958}.
Er leitet zudem Struktur her, die ihm nie vorgegeben wurde: Die Eingaben mit der höchsten Coverage
für cJSON und libxml2 sind tief verschachtelt, und der Anteil strukturell relevanter Bytes steigt
auf cJSON von \num{0.09} auf \num{0.90}.
Die Grenzen des Ansatzes liegen bei Durchsatz und Stabilität.
Ausführungsraten von \num{21} bis \num{23} Eingaben pro Sekunde im Gast liegen drei
Größenordnungen unter nativem coverage-gesteuertem Fuzzing, tiefere Suche zahlt sich jenseits von
50 Simulationen nicht mehr aus, und jedes Korpusziel endet unterhalb seines eigenen Höchstwerts
beim Anteil struktureller Bytes.
